\documentclass[10pt]{beamer}
\usetheme{metropolis}
\usepackage{graphicx}
\usepackage{listings}
\usepackage{booktabs}
\usepackage{xcolor}
\definecolor{codebg}{HTML}{F5F5F5}
\definecolor{codegreen}{HTML}{2E7D32}
\definecolor{codegray}{HTML}{757575}
\definecolor{codeblue}{HTML}{1565C0}
\lstdefinestyle{rstyle}{language=R,backgroundcolor=\color{codebg},basicstyle=\ttfamily\scriptsize,
  keywordstyle=\color{codeblue}\bfseries,stringstyle=\color{codegreen},
  commentstyle=\color{codegray}\itshape,breaklines=true,frame=single,
  rulecolor=\color{codegray!30},showstringspaces=false,tabsize=2,
  morekeywords={library,ggplot,aes,geom_line,geom_point,geom_segment,geom_col,geom_text,geom_bump,
    scale_color_manual,scale_y_reverse,theme_minimal,theme,element_blank,labs,
    fct_reorder,coord_flip,pivot_longer}}
\lstdefinestyle{pythonstyle}{language=Python,backgroundcolor=\color{codebg},basicstyle=\ttfamily\scriptsize,
  keywordstyle=\color{codeblue}\bfseries,stringstyle=\color{codegreen},
  commentstyle=\color{codegray}\itshape,breaklines=true,frame=single,
  rulecolor=\color{codegray!30},showstringspaces=false,tabsize=4,
  morekeywords={import,from,as,pd,plt,sns,np,fig,ax,plot,bar,text,invert_yaxis}}
\title{Rankings: Slope, Bump \& Waterfall}
\subtitle{Workshop 3 --- Module 6}
\author{Prof.Asoc.\ Endri Raco}
\institute{Data Visualization for Data Scientists \\ UNYT Tirana}
\date{2025/26}
\begin{document}

\begin{frame}
  \titlepage
\end{frame}

\begin{frame}{Module Roadmap}
  \begin{enumerate}
    \item Slopegraphs: two-point before/after comparison
    \item Bump charts: rank trajectories over multiple periods
    \item Waterfall charts: sequential cumulative breakdown
    \item Grey-accent strategy for slope and bump readability
    \item Chart selection: slope vs bump vs waterfall
  \end{enumerate}
  \begin{exampleblock}{Learning Outcome}
    You will produce slopegraphs for before/after comparisons, bump
    charts for rank trajectories, and waterfall charts for financial
    P\&L breakdowns --- three specialised chart types not covered by
    standard bar or line charts.
  \end{exampleblock}
\end{frame}

\section{Slopegraphs}

\begin{frame}{Slopegraph: Two-Point Comparison}
  \begin{center}
    \includegraphics[width=0.72\textwidth]{figures_w03m06/slopegraph}
  \end{center}
  \begin{exampleblock}{Data Insight}
    A slopegraph (Tufte, 1983) connects two values per item across two
    time points. Slope encodes rate of change; position encodes absolute
    value; crossing lines reveal rank changes. Use for 5--15 items across
    exactly two time points.
  \end{exampleblock}
\end{frame}

\begin{frame}{Slopegraph Best Practices}
  \begin{center}
    \includegraphics[width=0.92\textwidth]{figures_w03m06/slope_best_practices}
  \end{center}
  \begin{alertblock}{Pro-Tip}
    With $>$8 lines, slopegraphs become spaghetti. Apply the
    \textbf{grey-accent strategy}: draw all lines in light grey, then
    highlight 1--2 key items in colour with direct labels at both ends.
    This is the same technique used for scatter (W01-M04) and line charts.
  \end{alertblock}
\end{frame}

\begin{frame}[fragile]{Slopegraph in R vs Python}
  \begin{columns}[T]
    \begin{column}{0.48\textwidth}
      \textbf{R (ggplot2)}
\begin{lstlisting}[style=rstyle]
df_long <- df |>
  pivot_longer(c(v2020, v2024),
    names_to = "year",
    values_to = "value") |>
  mutate(year = parse_number(year))

ggplot(df_long, aes(
    x = year, y = value,
    group = category,
    color = highlight)) +
  geom_line(linewidth = 1) +
  geom_point(size = 3) +
  geom_text(data = . |>
    filter(year == 2020),
    aes(label = paste(category,
      scales::comma(value))),
    hjust = 1.1, size = 3) +
  geom_text(data = . |>
    filter(year == 2024),
    aes(label = paste(category,
      scales::comma(value))),
    hjust = -0.1, size = 3) +
  scale_color_manual(values =
    c("TRUE"="#E53935",
      "FALSE"="#CCCCCC")) +
  theme_void()
\end{lstlisting}
    \end{column}
    \begin{column}{0.48\textwidth}
      \textbf{Python}
\begin{lstlisting}[style=pythonstyle]
fig, ax = plt.subplots(
    figsize=(7, 5))

for cat, v0, v1, c in zip(
    cats, v2020, v2024, colors):
    # Grey-accent
    lw = 2.5 if cat == key else 1
    alpha = 1.0 if cat == key else 0.3

    ax.plot([0, 1], [v0, v1],
        "-o", color=c, lw=lw,
        markersize=6, alpha=alpha)

    # Direct labels at both ends
    ax.text(-0.08, v0,
        f"{cat} ({v0:,})",
        ha="right", fontsize=7,
        color=c)
    ax.text(1.08, v1,
        f"{cat} ({v1:,})",
        ha="left", fontsize=7,
        color=c)

ax.set_xticks([0, 1])
ax.set_xticklabels(["2020","2024"])
for sp in ax.spines.values():
    sp.set_visible(False)
ax.set_yticks([])
\end{lstlisting}
    \end{column}
  \end{columns}
\end{frame}

\section{Bump Charts}

\begin{frame}{Bump Chart: Rank Trajectories Over Time}
  \begin{center}
    \includegraphics[width=0.82\textwidth]{figures_w03m06/bump}
  \end{center}
\end{frame}

\begin{frame}[fragile]{Bump Chart in R vs Python}
  \begin{columns}[T]
    \begin{column}{0.48\textwidth}
      \textbf{R (ggbump)}
\begin{lstlisting}[style=rstyle]
# install.packages("ggbump")
library(ggbump)

ggplot(df_ranks, aes(
    x = year, y = rank,
    color = category)) +
  geom_bump(linewidth = 1.5,
    smooth = 8) +
  geom_point(size = 4) +
  geom_text(data = . |>
    filter(year == max(year)),
    aes(label = category),
    hjust = -0.2, size = 3,
    fontface = "bold") +
  scale_y_reverse() +
  scale_color_brewer(
    palette = "Set2") +
  theme_minimal() +
  theme(legend.position = "none",
    panel.grid.minor =
      element_blank()) +
  labs(title = "Category Rankings",
       x = "Year", y = "Rank")
\end{lstlisting}
    \end{column}
    \begin{column}{0.48\textwidth}
      \textbf{Python (manual)}
\begin{lstlisting}[style=pythonstyle]
fig, ax = plt.subplots(
    figsize=(9, 5))

for (cat, ranks), c in zip(
    rank_dict.items(), colors):
    # Grey-accent
    lw = 3 if cat in highlight else 1
    alpha = 1.0 if cat in highlight \
        else 0.3

    ax.plot(years, ranks, "-o",
        color=c, lw=lw, alpha=alpha,
        markersize=7)

    # Direct label at end
    ax.text(years[-1] + 0.2,
        ranks[-1], cat,
        fontsize=7, color=c,
        fontweight="bold",
        va="center")

ax.invert_yaxis()  # rank 1 = top
ax.set_yticks(range(1, n+1))
ax.set_yticklabels(
    [f"#{i}" for i in range(1,n+1)])
ax.grid(axis="y", alpha=0.1)
\end{lstlisting}
    \end{column}
  \end{columns}
  \begin{alertblock}{Pro-Tip}
    \texttt{scale\_y\_reverse()} / \texttt{ax.invert\_yaxis()} puts
    rank \#1 at the top. Direct-label at the rightmost point to
    eliminate the legend. Apply grey-accent when $>$4 categories.
  \end{alertblock}
\end{frame}

\section{Waterfall Charts}

\begin{frame}{Waterfall: Sequential Cumulative Breakdown}
  \begin{center}
    \includegraphics[width=0.82\textwidth]{figures_w03m06/waterfall}
  \end{center}
\end{frame}

\begin{frame}{P\&L Waterfall: Revenue to Net Income}
  \begin{center}
    \includegraphics[width=0.78\textwidth]{figures_w03m06/waterfall_pnl}
  \end{center}
  \begin{exampleblock}{Data Insight}
    Waterfall charts decompose a total into sequential additions
    (green) and subtractions (red). The connector lines show the
    running total. Standard in financial reporting (P\&L, budget
    variance). Start and end bars are blue (totals); intermediate
    bars are green (positive) or red (negative).
  \end{exampleblock}
\end{frame}

\begin{frame}[fragile]{Waterfall in R vs Python}
  \begin{columns}[T]
    \begin{column}{0.48\textwidth}
      \textbf{R (manual or waterfalls)}
\begin{lstlisting}[style=rstyle]
# Manual waterfall with geom_col
df <- df |>
  mutate(
    end = cumsum(value),
    start = lag(end, default = 0),
    color = case_when(
      type == "total" ~ "total",
      value > 0 ~ "positive",
      TRUE ~ "negative"))

ggplot(df, aes(
    x = fct_inorder(item))) +
  geom_col(aes(y = value,
    fill = color),
    width = 0.5) +
  geom_segment(aes(
    xend = lead(item),
    y = end, yend = end),
    linetype = "dashed",
    color = "grey50") +
  scale_fill_manual(values =
    c(total="#1565C0",
      positive="#2E7D32",
      negative="#E53935"))
\end{lstlisting}
    \end{column}
    \begin{column}{0.48\textwidth}
      \textbf{Python}
\begin{lstlisting}[style=pythonstyle]
for i in range(len(items)):
    if items[i] in ["Start","End"]:
        ax.bar(i, running[i],
            width=0.5,
            color="#1565C0",
            edgecolor="white")
    else:
        bottom = min(running[i-1],
                     running[i])
        height = abs(values[i])
        color = ("#2E7D32"
            if values[i] > 0
            else "#E53935")
        ax.bar(i, height,
            bottom=bottom,
            width=0.5,
            color=color,
            edgecolor="white")

    # Connector line
    if i < len(items) - 1:
        ax.plot([i+0.25, i+0.75],
            [running[i], running[i]],
            color="#888", lw=0.8,
            linestyle="--")
\end{lstlisting}
    \end{column}
  \end{columns}
\end{frame}

\section{Chart Selection}

\begin{frame}{Slope vs Bump vs Waterfall}
  \begin{center}
    \includegraphics[width=0.92\textwidth]{figures_w03m06/comparison_card}
  \end{center}
\end{frame}

\begin{frame}{When to Use Each}
  \centering
  \begin{tabular}{lp{4cm}l}
    \toprule
    \textbf{Scenario} & \textbf{Chart} & \textbf{Max Items} \\
    \midrule
    Before/after (2 points) & Slopegraph & 5--15 \\
    Rank over time (3--8 periods) & Bump chart & 4--10 \\
    Sequential breakdown to total & Waterfall & 5--12 steps \\
    Change comparison per category & Dumbbell (W03-M03) & 5--15 \\
    Simple trend over time & Line chart (W06) & 1--5 lines \\
    \bottomrule
  \end{tabular}
  \vspace{6pt}
  \begin{alertblock}{Pro-Tip}
    All three charts benefit from the \textbf{grey-accent strategy}:
    render most items in light grey, highlight 1--2 in colour. This
    prevents spaghetti and directs the reader to the finding.
  \end{alertblock}
\end{frame}

\section{Complete Examples}

\begin{frame}[fragile]{R: Slopegraph with Grey-Accent}
  \begin{columns}[T]
    \begin{column}{0.52\textwidth}
\begin{lstlisting}[style=rstyle]
df_long |>
  mutate(highlight =
    category == "FAMILY") |>
  ggplot(aes(x = year, y = value,
    group = category,
    color = highlight,
    alpha = highlight)) +
  geom_line(linewidth = 1.2) +
  geom_point(size = 3) +
  scale_color_manual(values =
    c("TRUE" = "#E53935",
      "FALSE" = "#CCCCCC")) +
  scale_alpha_manual(values =
    c("TRUE" = 1, "FALSE" = 0.4)) +
  # Direct labels at both ends
  geom_text(data = . |>
    filter(year == 2020),
    aes(label = paste(category,
      comma(value))),
    hjust = 1.1, size = 3) +
  theme_void() +
  theme(legend.position = "none") +
  labs(title = "FAMILY Surged
    to #1 by 2024")
\end{lstlisting}
    \end{column}
    \begin{column}{0.45\textwidth}
      \includegraphics[width=\textwidth]{figures_w03m06/r_output}
    \end{column}
  \end{columns}
\end{frame}

\begin{frame}[fragile]{Python: Bump Chart}
  \begin{columns}[T]
    \begin{column}{0.52\textwidth}
\begin{lstlisting}[style=pythonstyle]
fig, ax = plt.subplots(
    figsize=(9, 5))

for (cat, rk), c in zip(
    ranks.items(), colors):
    # Highlight top 2
    hl = cat in ["FAMILY", "GAME"]
    ax.plot(years, rk, "-o",
        color=c,
        lw=2.5 if hl else 1,
        alpha=1.0 if hl else 0.3,
        markersize=7 if hl else 4)
    ax.text(2024.2, rk[-1], cat,
        fontsize=7, color=c,
        fontweight="bold"
            if hl else "normal",
        va="center")

ax.invert_yaxis()
ax.set_yticks(range(1, 7))
ax.set_yticklabels(
    [f"#{i}" for i in range(1,7)])
ax.set_title(
    "FAMILY overtook GAME in 2021",
    fontweight="bold")
\end{lstlisting}
    \end{column}
    \begin{column}{0.45\textwidth}
      \includegraphics[width=\textwidth]{figures_w03m06/py_output}
    \end{column}
  \end{columns}
\end{frame}

\section{Synthesis}

\begin{frame}{Key Takeaways}
  \begin{enumerate}
    \item \textbf{Slopegraphs} show value change across exactly two time
          points. Slope = rate of change; crossing = rank reversal.
          5--15 items max.
    \item \textbf{Bump charts} track ordinal rank trajectories over 3--8
          periods. Invert the y-axis so \#1 is at the top. 4--10 items max.
    \item \textbf{Waterfall} charts decompose a total into sequential
          additions (green) and subtractions (red) with connector lines.
          Standard for P\&L and budget variance.
    \item All three charts require the \textbf{grey-accent strategy}
          to avoid spaghetti when items exceed 5--6.
    \item In R: manual \texttt{geom\_line + geom\_point + geom\_text},
          or \texttt{ggbump} for smooth bump curves.
          In Python: manual \texttt{ax.plot()} loops.
  \end{enumerate}
\end{frame}

\begin{frame}{Essential Readings}
  \begin{itemize}
    \item Tufte, E.~R. (1983). \emph{The Visual Display of Quantitative
          Information}, pp.\ 158--161 (Slopegraphs).
    \item Schwabish, J. (2021). \emph{Better Data Visualizations},
          Chapter 5 (Slope Charts, Bump Charts).
    \item Few, S. (2006). \emph{Information Dashboard Design},
          Chapter 6 (Waterfall charts in financial context).
    \item Wilke, C.~O. (2019). \emph{Fundamentals of Data Visualization},
          Chapter 11 (Visualizing Nested Proportions).
  \end{itemize}
  \vspace{6pt}
  \textbf{Next Module}: Uncertainty --- Error Bars, Confidence \& Gradient (W03-M07)
\end{frame}

\end{document}
